\documentclass[11pt,a4paper]{article}
\usepackage[T1]{fontenc}\usepackage[utf8]{inputenc}\usepackage{lmodern}
\usepackage[a4paper,margin=2.2cm]{geometry}
\usepackage{amsmath,amssymb}\usepackage{enumitem}\usepackage{booktabs,longtable,array}
\usepackage[hidelinks]{hyperref}
\setlength{\parindent}{0pt}\setlength{\parskip}{0.4em}
\title{\textbf{OP-grad F7: the $m_\sigma^2(u)$ branch --- can the completion
match the recorded $c_G=-1$?}\\\large Recon v1 --- proposal-only}
\author{Per reviewer directive after D4 acceptance (Variant A)}
\date{12 June 2026 --- rev.~1.1 (review wording safeguards applied)}
\begin{document}\maketitle

\section*{0. Scope and firewall}
G1 established $\gamma_{\kappa_n}=2-\gamma_{m_\sigma^2}-\gamma_{\hat a_{\rm
eff}}$ and showed that the recorded $c_G=-1$ is compatible with the kernel
route if the completion yields $m_\sigma^2\propto u$. This recon closes the
remaining bookkeeping tension at classification level: it assembles the
recorded constraint network on $m_\sigma$, shows that the $c_G$ exponent is a
\emph{sector-of-origin diagnostic} for the $\sigma$-gap, checks the linear
branch against every other recorded structure, and decomposes the derivation
target into the precise pair of completion functions. Proposal-only; no
preprint edits; nothing is claimed as derived microphysics.

\section{The recorded constraint network}
Two recorded facts already constrain $m_\sigma$ tightly.
\textbf{(a) Matching-point value.} The kernel chain
$\kappa_n\sim(\hat a_{\rm eff}/16\pi^2)\,u_0^2/m_\sigma^2\sim\phi_0^2$
(eq:MG\_u0) fixes the $\sigma$-gap at the parameter point:
\[
m_\sigma\;\simeq\;\sqrt{\frac{\hat a_{\rm eff}}{16\pi^2}}\;\frac{u_0}{\phi_0},
\]
heavy by construction (condensate-scale, up to the dimensionless cluster).
\textbf{(b) Slope.} The recorded $c_G=-1$, through G1 (with
$\gamma_{\hat a_{\rm eff}}=0$), fixes the local slope $\gamma_{m_\sigma^2}=1$.
Combining (a) and (b): the linear law
\[
m_\sigma^2(u)\;=\;c\,u,\qquad
c\;=\;\frac{\hat a_{\rm eff}}{16\pi^2}\,\chi,\qquad
\chi\;\equiv\;\frac{u_0}{\phi_0^2}\ \ (\text{dimensionless, since }[u]={\rm
GeV}^2),
\]
with a \emph{single} dimensionless cluster $(\hat a_{\rm eff},\chi)$ left for
the completion. F7 is therefore not an unconstrained function-hunting
problem: the value and the slope at the matching point are both fixed by the
record; only the \emph{origin} of the slope is open.

\section{The $c_G$ exponent as a sector-of-origin diagnostic}
The gradient sector has exactly one scale, the amplitude $u$ itself
($[u]={\rm GeV}^2$); the potential sector supplies the pinned scales
$\phi_0,\,m_\varphi^2\sim\lambda\phi_0^2$. Hence the possible origins of the
$\sigma$-gap map one-to-one onto the G1 branches:
\renewcommand{\arraystretch}{1.2}\small
\begin{longtable}{@{}p{4.3cm}p{2.6cm}p{1.5cm}p{4.6cm}@{}}
\toprule \textbf{Gap origin} & $m_\sigma^2$ & $c_G$ & \textbf{Status against
the record} \\ \midrule \endhead \bottomrule \endfoot
Gradient-native (no scale imported; only $u$ available) & $c\,u$
(dimensionally forced) & $-1$ & \emph{Recorded ansatz}; the linear law is the
unique dimensionally consistent gradient-native gap under the no-extra-scale
and fixed-$\hat a_{\rm eff}$ assumptions \\
Potential-imported (gap set by $\lambda\phi_0^2$) & pinned & $-2$ & Viable
alternative branch: requires closure-map renormalisation; all
$|\Delta\ln u|$-based numbers rescale by $1/2$ (thresholds relax
$\times2$); no new observable at current sensitivity \\
Mixed import ($u^2/\phi_0^2$) & $\propto u^2$ & $0$ & Incompatible with the
retained $\varepsilon$-sector interpretation: switches the recorded
$\varepsilon$-route off \\
\end{longtable}\normalsize
\textbf{Reformulation of F7.} The recorded $c_G=-1$, read through the kernel,
is precisely the statement that \emph{the heavy amplitude mode's gap is
generated within the gradient sector alone}. F7 thus reduces from ``derive
$m_\sigma^2(u)$'' to a binary structural question --- gradient-native versus
potential-imported gap --- with the third branch incompatible with the
retained $\varepsilon$-sector interpretation. Each answer has a definite,
bounded consequence; neither is an inconsistency.

\section{Consistency of the linear branch against the rest of the record}
\textbf{(a) No new light mode.} $m_\sigma\sim(\text{dimensionless
cluster})\times\phi_0$ stays condensate-scale on the entire cosmological
excursion ($|\Delta\ln m_\sigma^2|=|\Delta\ln u|\lesssim0.05$): the
heavy-$\sigma$ scale remains far above Hubble scales throughout the
excursion, so no additional rolling light degree of freedom is introduced.
\textbf{(b) No varying-constants leak.} $\sigma$ carries no recorded
holonomy or charge label (G2), so a $u$-tracking $m_\sigma$ feeds no gauge
threshold log: the protection route is untouched. The architecture is in
fact sharpened: \emph{$u$-tracking is mandatory in the neutral gravitational
channel ($\gamma_{m_\sigma^2}=1$ drives $c_G=-1$) and forbidden in the
charged channels ($|\gamma_{m_f^2}|\lesssim2\times10^{-2}$)} --- the same
charge-based selectivity, now acting on the mass spectrum.
\textbf{(c) $\varepsilon$-sector formulas unchanged.} With $c_G=-1$ the
recorded $u(z)=u_0(1+z)^{-2\varepsilon}$ and all derived bands stand as
written.

\section{The derivation target, decomposed}
What the completion must actually supply is not one function but a pair. The
canonical $\sigma$-gap is
\[
m_\sigma^2(u)\;=\;\frac{V''_{\rm eff}(u)}{Z_\sigma(u)},
\]
where $V_{\rm eff}(u)$ is the ordered-branch effective potential for the
amplitude and $Z_\sigma(u)$ the kinetic normalisation of its fluctuation
(both presently open --- they are the OP-grad core, the
$u_0\!\leftrightarrow\!(\mu^2,\lambda)$ relation in differential form). The
recorded $c_G=-1$ fixes one combination:
\[
\frac{\partial}{\partial\ln u}\,
\ln\!\frac{V''_{\rm eff}}{Z_\sigma}\;=\;1
\quad\text{at }u=u_0.
\]
A gradient-space symmetry-breaking sketch (Mexican-hat structure in the
gradient invariant, the same structure that selects the ordered branch)
naturally produces $V''_{\rm eff}$ and $Z_\sigma$ as powers of $u$ alone ---
the gradient-native scenario --- in which case the combination above is
dimensionally forced to an integer and the recorded value $1$ is the unique
choice compatible with a stabilised branch and a live $\varepsilon$-sector.
This sketch is a candidate route, not a derivation: the honest target for
the completion is the pair $(V_{\rm eff}(u),Z_\sigma(u))$ from the bare
$(\mu^2,\lambda,\alpha,\beta)$ data.

\section{Verdict}
\begin{enumerate}[leftmargin=1.9em,itemsep=2pt,label=\textbf{V\arabic*.}]
\item F7 is closed at classification level: the bookkeeping tension is
replaced by a binary structural question (gradient-native vs
potential-imported $\sigma$-gap), with the matching-point value and the
required slope both fixed by the record and the third branch incompatible
with the retained $\varepsilon$-sector interpretation.
\item The recorded $c_G=-1$ is equivalent, through the kernel, to the
gradient-native reading, in which the linear law $m_\sigma^2=c\,u$ is the
unique dimensionally consistent gap under the no-extra-scale and
fixed-$\hat a_{\rm eff}$ assumptions; the coefficient is fixed as
$c=(\hat a_{\rm eff}/16\pi^2)\,\chi$ with $\chi=u_0/\phi_0^2$.
\item The alternative $c_G=-2$ branch is viable, not falsified: it would
renormalise the closure map ($|\Delta\ln u|\to|\Delta\ln u|/2$, all
$\gamma$-thresholds relax $\times2$) with no new observable at current
sensitivity --- a bookkeeping decision deferred to the completion, exactly as
G1 framed it.
\item The linear branch passes every cross-check: no new light mode, no
varying-constants leak (charge selectivity now acts on the mass spectrum:
$u$-tracking mandatory in the neutral channel, forbidden in charged ones),
$\varepsilon$-sector formulas unchanged.
\item Named completion target: the pair $(V_{\rm eff}(u),Z_\sigma(u))$ from
bare data, with the recorded constraint
$\partial\ln(V''_{\rm eff}/Z_\sigma)/\partial\ln u=1$ at $u_0$. No preprint
edits from this front before review.
\end{enumerate}
\end{document}
